\documentclass[11pt]{article}
\usepackage{docstyle}
\renewcommand\sheetname{Homework \textperiodcentered{} Green Machines + Metals}
\begin{document}
\sheethead{Homework}{Green Machines + Metallic materials \textperiodcentered{} Year 9 Science}{about 80 minutes \textperiodcentered{} A / M / E}
Do Part 1 and Part 2 on two different days before the exam (about 40 minutes each).


\section{Part 1 \textperiodcentered{} Green Machines: Section A}
\Q Write the meaning of each word.\lvl{A}
\par\sub{a} germination
\alines{1}{when a root and shoot grow out of a seed and it develops into a seedling}
\sub{b} transpiration
\alines{1}{the loss of water vapour from a plant's leaves through the stomata}
\sub{c} zygote
\alines{1}{the cell formed when the pollen nucleus fuses with the ovule nucleus}

\Q Complete the sentences about flower parts.\lvl{A}
\par The \blank[26mm]{sepal} protects the flower while it is a bud. The \blank[26mm]{nectary} makes sweet nectar.
The \blank[26mm]{stigma} is sticky so that pollen grains stay on it.\par\vspace{2mm}

\Q Name three things a seedling competes with its parent plant for.\lvl{A}
\par 1. \blank[34mm]{light} \quad 2. \blank[34mm]{water} \quad 3. \blank[34mm]{minerals (or space)}\par\vspace{2mm}

\section{Part 1 \textperiodcentered{} Quick review}
\Q On a hot, dry afternoon a gardener sees the leaves of her tomato plants drooping.\lvl{A}
\par\sub{a} Which process is the plant losing water by? \hfill\blank[40mm]{transpiration}
\par\sub{b} Through which openings in the leaf is the water lost? \hfill\blank[40mm]{stomata}
\par\sub{c} Which tubes should be bringing water up from the roots? \hfill\blank[40mm]{xylem}\par\vspace{1mm}

\section{Part 1 \textperiodcentered{} Section B}
\Q Explain how a bee pollinates a flower. Use the names of the flower parts.\lvl{M}
\alines{4}{The bee is attracted by the brightly coloured petals, the scent and the nectar. Sticky pollen from the anthers rubs onto its body. When it visits another flower of the same species, the pollen rubs off onto the stigma. This is cross-pollination.}

\Q Explain why cross-pollination is better for a species than self-pollination.\lvl{E}
\alines{4}{Cross-pollination mixes the genetic material of two parents, so the offspring show more variation. If the environment changes or a disease arrives, some offspring are likely to have features that let them survive and reproduce, so the species continues.}

\Q Gardeners often soak bean seeds in water overnight before planting them. Explain why.\lvl{M}
\alines{3}{Water softens the testa, so the radicle can grow through it. Water is also one of the conditions for germination, so the seeds start to germinate sooner.}

\Q A carrot left in fresh water for a day becomes firm. A carrot left in very salty water becomes soft and bendy. Explain both results using osmosis.\lvl{E}
\alines{5}{Fresh water has a higher concentration of water than the carrot cells, so water moves into the cells by osmosis through the semi-permeable membranes and the carrot becomes firm. The salty water has a lower concentration of water than the cells, so water moves out of the cells by osmosis and the carrot goes soft.}

\Q A potted plant is weighed each day. All the mass it loses is water.\lvl{M}
\par\vspace{1mm}
\begin{tabular}{@{}l c c c c@{}}
\toprule
Day & 1 & 2 & 3 & 4\\
Mass (g) & 500 & 488 & 476 & 464\\
\bottomrule
\end{tabular}\par\vspace{1mm}
\sub{a} How much water does the plant lose each day?
\alines{1}{12 g per day ($500 - 488 = 12$ g)}
\sub{b} Explain where the water goes.
\alines{2}{It evaporates from the leaves and is lost as water vapour through the stomata (transpiration).}

\Q Sam explains how water gets to the leaves of a plant.\lvl{E}
\par\vspace{1mm}
\samcard{0.9\linewidth}{\flawed{Water enters the root hairs by osmosis, from a low concentration of water to a high concentration of water. It travels up the phloem to the leaves. The leaves make glucose and oxygen from carbon dioxide and water. At night the plant stops respiring.}}\par\vspace{1mm}
\textbf{Sam's answer has 3 errors.} Find each one, fix it, and say why it loses marks.
\alines{5}{1. Osmosis is from a high to a low concentration of water (\Lref{osmosis}). 2. Water travels up the xylem, not the phloem (\Lref{xylem}). 3. Plants respire all the time, day and night (\Lref{respire}).}

\vspace{4mm}
{\small Also set: Green Machines booklet p.24, fill in the glossary for any ten words.}

\clearpage
\section{Part 2 \textperiodcentered{} Metallic materials: Section A}
\Q Write the meaning of each word.\lvl{A}
\par\sub{a} density
\alines{1}{how much mass is in a given volume (mass divided by volume)}
\sub{b} alloy
\alines{1}{a mixture of two or more elements, at least one of which is a metal}
\sub{c} corrosion
\alines{1}{a chemical reaction of a metal with substances around it, such as oxygen and water}

\Q Name the heat treatment.\lvl{A}
\par\sub{a} A blacksmith wants a horseshoe to be soft, so it is easy to shape. \hfill\blank[34mm]{annealing}
\par\sub{b} The edge of a chisel must be very hard. \hfill\blank[34mm]{quenching}\par\vspace{1mm}

\Q Write the word equations.\lvl{A}
\par\sub{a} lead + oxygen $\longrightarrow$ \blank[70mm]{lead oxide}
\par\sub{b} sodium + hydrochloric acid $\longrightarrow$ \blank[70mm]{sodium chloride + hydrogen}\par\vspace{1mm}

\section{Part 2 \textperiodcentered{} Quick review}
\Q Name the property of metals that makes each one a good choice.\lvl{A}
\par\sub{a} the copper base of a saucepan \hfill\blank[50mm]{good conductor of heat}
\par\sub{b} gold contacts inside a phone \hfill\blank[50mm]{conducts electricity, unreactive}
\par\sub{c} a steel bridge cable \hfill\blank[50mm]{strong}\par\vspace{1mm}

\section{Part 2 \textperiodcentered{} Section B}
\Q A gold-coloured crown has a mass of 965\,g and a volume of 50.0\,cm$^{3}$. The density of pure gold is 19.3\,g/cm$^{3}$. Is the crown pure gold? Show your working.\lvl{M}
\flines{2}{density $= \dfrac{965\ \mathrm{g}}{50.0\ \mathrm{cm^{3}}} = 19.3$ g/cm$^{3}$\quad This matches pure gold, so it could be pure gold.}

\Q Aircraft bodies are made of an aluminium alloy, not steel. Use the data to explain why.\lvl{E}
\par\vspace{1mm}
\begin{tabular}{@{}l c c@{}}
\toprule
Material & Density (g/cm$^{3}$) & Strength\\
\midrule
aluminium alloy & 2.8 & high\\
steel & 7.8 & very high\\
\bottomrule
\end{tabular}
\alines{4}{The aluminium alloy is strong enough, and its density (2.8) is much lower than steel's (7.8), almost three times less. A lighter aircraft needs less fuel to fly. Steel is stronger, but it would make the aircraft far too heavy.}

\Q The steel exhaust pipe of a car rusts faster than the steel body of the car. Explain why.\lvl{M}
\alines{3}{The exhaust pipe gets hot, and the burning petrol makes water, so the pipe has water, oxygen and a high temperature. A raised temperature speeds up rusting.}

\Q A student adds the same mass of four metals to the same volume of dilute acid and measures the temperature rise.\lvl{E}
\par\vspace{1mm}
\begin{tabular}{@{}l c c c c@{}}
\toprule
Metal & magnesium & zinc & iron & copper\\
Temperature rise (\textdegree C) & 18 & 7 & 3 & 0\\
\bottomrule
\end{tabular}\par\vspace{1mm}
\sub{a} Put the metals in order of reactivity, most reactive first.
\alines{1}{magnesium, zinc, iron, copper}
\sub{b} Use the data to explain your order, and name two things the student kept the same.
\alines{4}{The more reactive the metal, the faster it reacts with the acid, so it gives out more heat: magnesium raises the temperature most (18\,\textdegree C) and copper does not react at all (0\,\textdegree C). Kept the same: mass of metal, volume of acid (also concentration of acid, starting temperature).}

\Q Sam writes about alloys and reactions.\lvl{E}
\par\vspace{1mm}
\samcard{0.9\linewidth}{\flawed{1. Steel is harder than pure iron because all the atoms in steel are the same size, so the layers slide easily.\\
2. magnesium + hydrochloric acid makes magnesium chloride + water\\
3. Bronze is an alloy, so it is a mixture of metals.\\
4. When copper is added to dilute hydrochloric acid, bubbles of hydrogen are given off.}}\par\vspace{1mm}
\textbf{Sam's answer has 3 errors.} Find each one, fix it, and say why it loses marks.
\alines{5}{1. Steel has atoms of different sizes (carbon), which disrupt the layers so they cannot slide easily (\Lref{alloy}). 2. Metal + acid gives hydrogen, not water: magnesium chloride + hydrogen (\Lref{wordeq}). 3. Copper is below the acids in the reactivity series, so there are no bubbles: no reaction (\Lref{react}).}

\vspace{4mm}
{\small Also set: Metals booklet p.10--11 (draw the solder graph, then answer Q1--2) and the crossword on p.16.}
\end{document}
